Referee machine
ESP32 vs Python Server
Option 1 — Referee on the ESP32 master (C++ firmware)
The master owns the game, drives the LEDs and feeder directly, and serves a thin web page to the visitor's phone over the WiFi access point it already runs.
Pros:
•	It's mostly already built. loopGameModeM4 is a working referee; you'd extend rather than start over. For a time-boxed thesis, reusing working code is a real advantage.
•	It's one self-contained box. No extra computer to power, cool, secure, or reboot — fewer things that can fail in an enclosure no one is watching at night. Embedded boards are designed to run unattended for months and boot instantly on power-up.
•	The referee sits right next to the hardware it controls. "Decide match → light the LED" has no network hop, so the chimp's LEDs and the game state can't drift apart.
•	Cheaper, lower power, less to maintain. The existing topology already centers on the master (the feeder is already a client to it), so this fits with the least disruption.
Cons:
•	It's C++ on a microcontroller — the unforgiving environment you've seen (the uint8_t overflow, Serial-print debugging, recompile-and-flash for every change). Slower and more error-prone for you specifically.
•	The ESP32 is small. Serving a web app from it is doable but you feel the ceiling — limited memory, fiddlier WebSocket handling, no room for rich extras.
•	Every change means flashing the board. Iteration is slower than editing a file and refreshing a browser.
Option 2 — Referee on a Python server (a separate computer, realistically a Raspberry Pi)
A Python program (something like Flask/FastAPI) runs the referee and serves the web UI; the ESP boards become "dumb I/O" that just read sensors/buttons and light LEDs when the server tells them to.
Pros:
•	You write in Python — comfortable, fast, easy to debug, easy to change. You'd likely reach "working" faster despite the system being bigger.
•	A real computer can serve a polished web app, handle flaky phone connections gracefully, and host any future extras (dashboards, attract-mode assets, remote monitoring) without fighting for memory.
•	Edit-and-restart iteration, no flashing.
Cons:
•	You're adding a whole computer to an outdoor public installation — another thing to power, cool, secure, update, and that can crash or corrupt its SD card. That's a genuine reliability and maintenance liability for an unattended kiosk, and the brief already flags summer cooling as a concern.
•	The referee is now a network hop away from the hardware, so "light this LED" crosses WiFi/serial — more latency and more ways for the phone view and the LED view to fall out of sync.
•	You don't escape firmware. The ESPs still need code to read sensors and drive LEDs; you've just made that code dumber and added a new server-to-ESP protocol to design. So it's often more total work, across two codebases that must agree, and you throw away the working referee the master already has.
Alternatives worth naming
•	Referee in the browser (your phone runs the logic in JavaScript). This is closest to your current prototype and the easiest web development. But it's the wrong fit for the real system: the chimp has no phone, the chimp's LEDs would depend on a specific visitor's phone staying connected, and if they walk away mid-game the state is gone. It's perfect for a standalone visitor toy or a demo, and a poor choice once two displays must agree. Good to keep in your back pocket for the no-hardware version, not for the installation.
•	Raspberry Pi specifically is just the sensible embodiment of Option 2 — small, cheap, low-power enough to be less of a packaging headache than a full PC, while still being a Linux box you have to keep alive (SD-card corruption is the classic Pi failure). If you go the Python-server route, this is what it'd actually be.
 
 
First, the core idea you're reaching for is real and has a name: multiplexing. One sensing channel, shared across many pads, by either asking "which pad?" in sequence or wiring the pads in a grid so a position is identified by a row-plus-column combination. So yes — one sensor circuit reading 16 pads is absolutely a thing people do. But notice I said one sensor circuit, not one KI6001. That distinction is the whole answer, so let me walk through the real options from "swap the expensive part" to "one channel for everything."
Option A — keep the one-per-field design but use a far cheaper sensor. Your KI6001 is an industrial proximity sensor — rugged, certified, and priced for factory automation, which is why your brief calls it expensive. But "detect a hand through glass" doesn't require an industrial part. A capacitive touch controller chip does the same job for a tiny fraction of the cost. The relevant idea is self-capacitance sensing: a conductive pad behind the glass holds a small electrical charge, your hand (being conductive and grounded-ish) changes that charge when it approaches, and a chip measures the change. Cheap chips built exactly for this exist — and the ESP32 you already use has built-in capacitive touch inputs on many of its pins, meaning the "sensor" can be a copper pad plus firmware, with no dedicated sensor part at all. So the first, lowest-risk win isn't "fewer sensors," it's "same architecture, each field senses with a copper pad + the ESP's own touch peripheral (or a ~1 € touch chip) instead of a costly industrial unit." This alone could be most of your cost saving, and it changes nothing about your game logic.
The catch with that, and it's important: glass thickness. Capacitive sensing through a barrier gets dramatically harder as the barrier thickens, because the effect falls off fast with distance. Your security glass is thick, and that's precisely why the previous person chose an industrial sensor with the power to punch through it. So before you trust any cheaper option, the very first thing to do — before design, before buying in bulk — is take one candidate sensor and one piece of the actual enclosure glass and measure whether a hand reliably registers through it. Everything downstream depends on that one bench test. Don't skip it; it's the make-or-break.
Option B — multiplex several pads onto one sensing chip. Here's where "1 sensor, 16 pads" becomes literal. Touch-controller chips are commonly multi-channel — a single chip might read 8, 12, or 16 pads at once, internally scanning them and telling you over a data line which pad was touched. So "16 fields, one sensing chip" is an off-the-shelf reality, not a hack. The cost: you run a wire from every pad back to wherever that chip lives, and capacitive sensing hates long wires — the wire itself acts like an antenna and adds stray capacitance that drowns out the small change a finger makes. So a central chip with 16 long pad-wires fighting electrical noise is exactly the kind of reliability headache you're already facing with I²C. This option trades sensor cost for wiring difficulty, which in your installation is the scarce resource. It can work, but it pushes hard on the thing that's hardest in your build.
Option C — a true matrix (the "one channel, 16 positions" dream). Arrange the 16 pads as a 4×4 grid: 4 row lines, 4 column lines, and a touch is identified by which row and which column light up together. This is how keypads and touchscreens get away with few wires for many keys — 8 lines instead of 16, a position from the crossing. Elegant on paper. But for your game it has a specific, disqualifying-ish problem, and this is the part most people miss: a matrix struggles when two pads in the same row or column are touched at once — it can't always tell a real two-touch from a "ghost" crossing. Standard keypads assume one key at a time. Now look at your game: it's a two-player competitive game where both players can be reaching at the board, and the whole mechanic is comparing two positions. The moment you can't trust simultaneous or rapid two-touches, the matrix's cheapness has bought you a bug in the exact heart of your game. So the most elegant wiring solution is the worst fit for this particular application, and that's worth knowing before you fall in love with it.
So, putting it together for your situation, the honest hierarchy: the biggest, safest saving is Option A — replace the costly industrial sensor with a cheap capacitive-touch approach per field (copper pad + the ESP's built-in touch, or a low-cost touch chip), keeping one sensing channel per field so your game logic and reliability profile don't change. Multiplexing further (Option B) is genuinely possible and is how you'd reach "one chip for many pads," but it spends your scarcest resource — clean wiring — and only do it if the per-field-chip cost still bothers you after Option A. The full matrix (Option C) I'd avoid for this game specifically, because two-player simultaneous touch is a feature you can't lose.
And every one of these lives or dies on the same gate: does it sense a hand through your real glass? That single measurement, done early on a scrap of the actual glass, tells you which of these options are even on the table. A cheap sensor that can't see through 15mm of security glass saves you nothing.
One reframe to leave you with, because it reorders your shopping: your task says "least/cheaper sensors," but the constraint is the glass, and the cost is split between the sensor part and the wiring it demands. The cheapest robust answer is usually "cheap per-field sensing, short wires" — not "one heroic sensor doing everything." Fewer sensors sounds cheaper, but if it forces long noisy wires or a matrix that ghosts on two-touch, you pay it back in debugging and unreliability, which in a thesis is the most expensive currency you have.
 
 
Let me build you a clean one — positioning your own project as the final row, which is what makes it work for a related work section.Here it is — the four examples plus your own project as the final row, so the table does double duty as positioning for your contribution.
System	Animal / site	Interaction & barrier	Reward	Game / task	Relevance to your project
Interactive projection (Zoos Victoria + Univ. Melbourne)	Orangutans, Melbourne Zoo	Body/hand tracked by a Kinect depth camera; projector + sensor sit outside the enclosure, no contact hardware	None — deliberately reward-free	Simple play: "catch the dot," chase, music, painting; trialled human–ape play across the glass	Same "all tech outside the barrier" goal, solved optically; reward-free led to brief, intermittent use
Apex + ApeTouch (Indianapolis Zoo / Kyoto; commercial)	Macaques & great apes, multiple zoos	Touchscreen at the cage mesh; portable, attach/detach with animals present; separate human-facing control screen	Food + tone per correct touch	Task suite incl. Matching Pairs (Memory), image choice, numerical sequencing; difficulty levels + play-vs-computer	"Buy vs build" benchmark; ready-made Memory game and difficulty modes; its human screen ≈ your planned web UI
Gorilla Game Lab (Univ. Bristol / Bristol Zoo; custom)	Western lowland gorillas	Physical poke-holes worked with fingers/tools; screen-free; hidden sensors + cameras auto-log use	Food (nuts) drawn through the modules	Reconfigurable puzzle: 12 interchangeable modules → a new layout every trial (anti-habituation)	Closest to your custom build; modularity for novelty; cheap hidden logging; a voluntary-use welfare-test template
Orangutan Tic-Tac-Toe / Memory (Indianapolis Zoo)	Orangutans	Touchscreen at a protected, visitor-facing interface	Food (typical for these tasks)	Familiar competitive games: Tic-Tac-Toe and Matching Pairs/Memory	Precedent that a recognisable game (Memory) is simultaneously enrichment and visitor education
Memoryautomat — this work (Zoo Karlsruhe)	Chimpanzees	Capacitive touch through security glass (chimp side); buttons/web UI (human side); all tech outside the enclosure	Food via a networked feeder on a win	Colour-based Memory (matching pairs); human-vs-(eventually-)chimp; scaling 4 → 16 fields	Sits between a touchscreen and a tactile puzzle: a physical, food-rewarded, two-player-across-glass device
For the reference list behind it: the Melbourne projection is Carter, Sherwen & Webber (2021, Zoo Biology) plus Webber et al. (2017, CHI); Apex/ApeTouch is Martin, Muramatsu & Matsuzawa (2022, Animals); and the Gorilla Game Lab is Clark et al. (2019, Frontiers in Psychology) with the companion Gray et al. (2018, ACI). The Indianapolis games are best cited via the ApeTouch paper and zoo/popular coverage rather than a dedicated study.
A couple of honest caveats so you don't overstate it: the Indianapolis row is the thinnest — it's well documented in zoo/press material but not in a single dedicated paper, so keep its claims modest. And the reward column is a real axis of disagreement in this literature (reward-free vs food rewarded), so it's worth a sentence in your text rather than letting the table imply everyone agrees.
Want me to drop this into a Word or LaTeX table so it goes straight into your thesis, or expand it into a couple of paragraphs of related work prose that walks through the table?
 
 
 

